To assess the usefulness of the framework beyond the next hour, the rolling-origin protocol was repeated with direct forecasting at $h \in \{1, 6, 12, 24, 48\}$ hours (Table~\ref{tab_horizons}, Figure~\ref{fig_horizons}). A separate model was trained for each horizon, and the baselines used the same information available at the forecast origin. At ATE, the MAE of the best model rose from 3.12 (HGB, $h=1$) to 5.63\,$\mu$g/m$^3$ (HGB, $h=48$), an increase of 80\%; it outperformed the naive forecast at $h \in \{1, 6, 12, 24, 48\}$ and the 24-h seasonal naive at $h \in \{1, 6, 12, 24, 48\}$. At Campo de Marte, the MAE of the best model rose from 6.10 (RF, $h=1$) to 14.28\,$\mu$g/m$^3$ (RF, $h=48$), an increase of 134\%; it outperformed the naive forecast at $h \in \{1, 6, 12\}$ and the 24-h seasonal naive at $h \in \{1, 6, 12\}$. At San Borja, the MAE of the best model rose from 3.22 (MLP, $h=1$) to 6.86\,$\mu$g/m$^3$ (HGB, $h=48$), an increase of 113\%; it outperformed the naive forecast at $h \in \{1, 6, 12\}$ and the 24-h seasonal naive at $h \in \{1\}$. At Santa Anita, the MAE of the best model rose from 3.36 (MLP+met, $h=1$) to 7.55\,$\mu$g/m$^3$ (SARIMA, $h=48$), an increase of 125\%; it outperformed the naive forecast at $h \in \{1, 6, 12, 48\}$ and the 24-h seasonal naive at $h \in \{1, 48\}$. For horizons of 12 to 48 h, the best model outperformed the 24-h seasonal naive in 5 of 12 station--horizon combinations (42\%). This result delimits the operational scope of the univariate formulation: its advantage is concentrated in the short term, whereas at daily horizons the persistent diurnal profile captures much of the predictable information.
